
\subsubsection{6.1 Stratum Collapse}

The central finding of this experiment is that geometry-based contamination routing has a hard design prerequisite that was not implemented: top-k nearest-neighbor corpus retrieval for SBERT cosine computation. Any implementation that substitutes random corpus sampling for top-k retrieval when computing embedding-based similarity features is expected to produce degenerate score distributions at production scale. Score distributions change qualitatively, not merely quantitatively, as corpus sample size grows under random sampling — a property that dry-run validation at small scales does not reveal.

The stratum collapse phenomenon has not been previously described in the contamination detection literature. The conditions that produce it — random corpus subsampling for embedding similarity features — are present in some prior work, suggesting that semantic stratum formation may be underestimated in earlier studies using random corpus subsets.

\subsubsection{6.2 Benchmark Heterogeneity}

The finding that MMLU achieves n-gram recall = 1.0 and GSM8K achieves recall = 0.0 against the same corpora reveals that heterogeneous benchmark sets have structurally different corpus overlap profiles. A single recall criterion applied uniformly across benchmark types averaging over NLU and math benchmarks produces a statistic that is not diagnostic of either type. Future contamination detection frameworks should apply benchmark-type-specific criteria and evaluate math benchmarks against domain-specific corpora (e.g., OpenWebMath, Proof-Pile) rather than general web corpora.

\subsubsection{6.3 Hypothesis Status}

The MUST_WORK gate failure for h-e1 is attributed to experimental design and implementation failures:

\begin{enumerate}
\item Random corpus streaming for SBERT cosine computation (design flaw)
\item DCPDDDetector log-ratio bug (engineering bug)
\item ConStatDetector non-standard API usage (engineering bug)
\item NgramIndex total count vs. max consecutive run (implementation deviation)
\item StratifiedEvaluator scalar broadcast (engineering bug)
\item FineWeb substitution for RedPajama (data substitution)
\item Heterogeneous benchmark types in unified recall criterion (criterion mismatch)
\end{enumerate}

None of these constitute falsifying evidence against the hypothesis. The theoretical causal chain — that detector families operate on orthogonal signal types whose efficacy is structurally determined by corpus overlap geometry — is consistent with the contamination detection literature [Fu et al., 2024; Singh et al., 2024] and has not been tested. Hypothesis confidence was revised from 0.78 to 0.62, reflecting the severity of the design flaw, not a negative empirical finding.

Dependent sub-hypotheses h-m1 through h-m4 are cascade-failed pending redesign of h-e1.

\subsubsection{6.4 Limitations}

\textbf{L1 (Critical):} Stratum collapse invalidates the primary routing accuracy measurement (P1) and indeterminacy rate measurement (P3). The semantic stratum contains zero items; all stratum-conditional metrics are undefined or vacuous.

\textbf{L2 (Moderate):} Implementation bugs B3 and B4 make all MIA-based and statistical detector results uninterpretable. Min-K\%++ F1 variance = 0.000 reflects Bug B3, not true detector behavior.

\textbf{L3 (Moderate):} Corpus coverage at 50K documents represents < 6\% of The Pile by document count. All n-gram recall figures are lower bounds on true corpus overlap. GSM8K recall = 0.0 may be a sampling artifact.

\textbf{L4 (Low-Moderate):} FineWeb was substituted for RedPajama-Data-1T without prior specification. FineWeb has different deduplication characteristics. Cross-corpus generalization claims involving FineWeb require re-evaluation with RedPajama.

\textbf{L5 (Low):} Dry-run results (3-way stratification at 500-document scale) were not predictive of full-scale behavior. Small-scale dry runs are insufficient for validating SBERT-based stratum formation.

\textbf{L6 (Low):} The stratum collapse boundary condition is documented for \texttt{all-MiniLM-L6-v2}. Whether it generalizes to other SBERT model variants is unverified.

